% TODO make hw stylesheet
\documentclass[10pt]{article}
\usepackage[letterpaper,top=1in,left=1in,right=1in]{geometry}
\usepackage[dvipsnames,svgnames]{xcolor}
\usepackage[shortlabels]{enumitem}
\usepackage[framemethod=TikZ]{mdframed}
\usepackage{amsmath,amssymb,amsthm}
\usepackage[colorlinks]{hyperref}
\usepackage{multicol}
\usepackage{tabularx}
\usepackage{bbm}

\hypersetup{
    colorlinks=true,
    linkcolor=blue,
    filecolor=magenta,
    urlcolor=magenta,
    pdftitle={MAT 324 HW}
}

\usepackage{mathtools}
\usepackage{thmtools}
\usepackage[textsize=footnotesize]{todonotes}
\usepackage{titling}
\usepackage{cancel}

\reversemarginpar
\mdfdefinestyle{mdbluebox}{roundcorner=10pt,innerbottommargin=9pt,
linecolor=blue,backgroundcolor=TealBlue!5,}
\declaretheoremstyle[headfont=\sffamily\bfseries\color{MidnightBlue},
mdframed={style=mdbluebox},]{thmbluebox}

\mdfdefinestyle{mdbloobox}{frametitlefont=\bfseries,innerbottommargin=8pt,
nobreak=true,backgroundcolor=TealBlue!5,linecolor=blue,}
\declaretheoremstyle[headfont=\bfseries\color{MidnightBlue},
mdframed={style=mdbloobox},headpunct={\\[3pt]},postheadspace=0pt,]{thmbloobox}

\mdfdefinestyle{mdredbox}{frametitlefont=\bfseries,innerbottommargin=8pt,
nobreak=true,backgroundcolor=Salmon!5,linecolor=RawSienna,}
\declaretheoremstyle[headfont=\bfseries\color{RawSienna},
mdframed={style=mdredbox},headpunct={\\[3pt]},postheadspace=0pt,]{thmredbox}

\mdfdefinestyle{mdgreenbox}{linecolor=ForestGreen,backgroundcolor=ForestGreen!5,
linewidth=2pt,rightline=false,leftline=true,topline=false,bottomline=false,}
\declaretheoremstyle[headfont=\bfseries\sffamily\color{ForestGreen!70!black},
mdframed={style=mdgreenbox},headpunct={ --- },]{thmgreenbox}

\mdfdefinestyle{mdorangebox}{linecolor=RedOrange,backgroundcolor=RedOrange!5,
linewidth=2pt,rightline=false,leftline=true,topline=false,bottomline=false,}
\declaretheoremstyle[headfont=\bfseries\sffamily\color{RedOrange!70!black},
mdframed={style=mdorangebox},headpunct={ --- },]{thmorangebox}

\mdfdefinestyle{mdblackbox}{linecolor=black,backgroundcolor=RedViolet!5!gray!5,
linewidth=3pt,rightline=false,leftline=true,topline=false,bottomline=false,}
\declaretheoremstyle[mdframed={style=mdblackbox}]{thmblackbox}

\declaretheoremstyle[spaceabove=6pt,spacebelow=6pt]{basehead}

\declaretheorem[style=basehead,name=Problem]{problem}
\declaretheorem[style=thmredbox,name=Problem,numbered=no]{problem*}
\declaretheorem[style=thmbloobox,name=Setup]{setup} % for config geo
\declaretheorem[style=thmredbox,name=Example,sibling=problem]{example}
\declaretheorem[style=thmredbox,name=$\bigstar$ Problem,sibling=problem]{niceprob}
\declaretheorem[style=thmbluebox,name=Theorem,sibling=problem]{theorem}
\declaretheorem[style=thmbluebox,name=Corollary,sibling=theorem]{corollary}
\declaretheorem[style=thmbluebox,name=Proposition,sibling=theorem]{prop}
\declaretheorem[style=thmbluebox,name=Lemma,numberwithin=problem]{lemma}
\declaretheorem[style=thmbluebox,name=Theorem,numbered=no]{theorem*}
\declaretheorem[style=thmbluebox,name=Lemma,numbered=no]{lemma*}
\declaretheorem[style=thmgreenbox,name=Claim,numberwithin=problem]{claim}
\declaretheorem[style=thmgreenbox,name=Claim,numbered=no]{claim*}
\declaretheorem[style=thmblackbox,name=Remark,numberwithin=problem]{remark}
\declaretheorem[style=thmblackbox,name=Remark,numbered=no]{remark*}
\declaretheorem[style=thmgreenbox,name=Definition,sibling=theorem]{definition}
\declaretheorem[style=thmgreenbox,name=Definition,numbered=no]{definition*}

\providecommand{\ol}{\overline}
\providecommand{\ceil}[1]{\left\lceil #1 \right\rceil}
\providecommand{\floor}[1]{\left\lfloor #1 \right\rfloor}
\newcommand{\dd}{\,\mathrm{d}}
\providecommand{\ul}{\underline}
\providecommand{\eps}{\varepsilon}
\providecommand{\half}{\frac{1}{2}}
\providecommand{\CC}{\mathbb C}
\providecommand{\FF}{\mathbb F}
\providecommand{\NN}{\mathbb N}
\providecommand{\QQ}{\mathbb Q}
\providecommand{\RR}{\mathbb R}
\providecommand{\ZZ}{\mathbb Z}
\providecommand{\ind}{\mathbbm 1}
\providecommand{\dg}{^\circ}
\providecommand{\ii}{\item}
\providecommand{\setdiff}{\smallsetminus}
\providecommand{\alert}[1]{{\sffamily\textbf{\textcolor{blue}{#1}}}}
\providecommand{\inv}{^{-1}}
\providecommand{\mb}[1]{\mathbf{#1}}

\newenvironment{soln}{\begin{proof}[Solution]}{\end{proof}}

\providecommand{\pow}{\operatorname{Pow}}
\providecommand{\vocab}[1]{{\textbf{\textcolor{ForestGreen}{#1}}}}
\providecommand{\lcm}{\operatorname{lcm}}
\providecommand{\abs}[1]{\left\lvert #1\right\rvert}
\providecommand{\norm}[1]{\left\| #1\right\|}
\providecommand{\gr}{\operatorname{gr}}

\usepackage{mathrsfs}

% place title at top right
\setlength{\droptitle}{-8ex}
\pretitle{\begin{flushright}\Large\bfseries}
\posttitle{\par\end{flushright}}
\preauthor{\begin{flushright}}
\postauthor{\end{flushright}}
\predate{\begin{flushright}}
\postdate{\end{flushright}}

\title{MAT 324 HW07}
\author{\large Aditya Pahuja}
\date{\large Due 10/30}
\begin{document}
\maketitle

\begin{problem}
    Suppose $X$, $Y$ are sets, $\mathcal S\subseteq 2^X$ with $X\in\mathcal S$,
    and $\mathcal T\subseteq 2^Y$ with $Y\in\mathcal T$.
    \begin{enumerate}[(a)]
        \ii Show that $\sigma(\mathcal S)\otimes\sigma(\mathcal T) =
	\sigma(\mathcal S\times\mathcal T)$, where
	$\{A\times B : A\in\mathcal S, B\in\mathcal T\}$.
	\ii Given an example showing that neither assumption $X\in\mathcal S$
	nor $Y\in\mathcal T$ can be dropped for the conclusion above.
    \end{enumerate}
\end{problem}

\begin{soln}
    (a) On one hand $\sigma(\mathcal S\times\mathcal T)\subseteq
    \sigma(\mathcal S)\otimes\sigma(\mathcal T)$ because
    the latter set contains
    $\sigma(\mathcal S)\times\sigma(\mathcal T)
    \supseteq \mathcal S\times\mathcal T$.

    To show the reverse inclusion, we prove
    $A\times B\in \sigma(\mathcal S\times\mathcal T)$ for all
    $A\in\sigma(\mathcal S)$, $B\in\sigma(\mathcal T)$.
    Given $A\in\mathcal S$, let
    \begin{equation*}
	\mathcal M_A
	= \{B\in\sigma(\mathcal T) : A\times B\in \sigma(\mathcal S\times\mathcal T)\}.
    \end{equation*}
    Then $\mathcal M_A$ is a $\sigma$-algebra on $Y$:
    \begin{enumerate}[(1)]
        \ii $\varnothing\in\mathcal M_A$, since
	$A\times\varnothing = \varnothing\in \sigma(\mathcal S\times\mathcal T)$.
	\ii Given $B_1,B_2,\ldots\in\mathcal M_A$,
	$\bigcup_{n=1}^\infty B_n\in\mathcal M_A$ because
	\begin{equation*}
	    A\times\bigcup_{n=1}^\infty B_n
	    = \bigcup_{n=1}^\infty \underbrace{A\times B_n}_
	    {\in \sigma(\mathcal S\times\mathcal T)}
	    \in \sigma(\mathcal S\times\mathcal T).
	\end{equation*}
	\ii Given $B\in\mathcal M_A$,
	the complement $B^c$ is in $\mathcal M_A$ because
	\begin{equation*}
	    A\times B^c = A\times Y\setminus A\times B\in
	    \sigma(\mathcal S\times\mathcal T),
	\end{equation*}
	since $A\times Y\in\mathcal S\times\mathcal T$
	because $Y\in\mathcal T$ and $A\times B\in\sigma(\mathcal S\times\mathcal T)$
	by assumption.
    \end{enumerate}
    Additionally, $\mathcal T\subseteq\mathcal M_A$ since
    $\{A\}\times\mathcal T\subseteq\sigma(\mathcal S\times\mathcal T)$,
    so $\mathcal M_A = \sigma(\mathcal T)$.
    Now, given $B\in \sigma(\mathcal T)$, let
    \begin{equation*}
	\mathcal M^B =
	\{A\in\sigma(\mathcal S) : A\times B\in \sigma(\mathcal S\times\mathcal T)\}.
    \end{equation*}
    This is checked to be a $\sigma$-algebra on $X$
    in the exact same way that we checked that $\mathcal M_A$ is a $\sigma$-algebra.
    Moreover, given $A\in\mathcal S$ and $B\in\sigma(\mathcal T) = \mathcal M_A$,
    we have $A\times B \in\sigma(\mathcal S\times\mathcal T)$,
    which implies that $A\in\mathcal M^B$ as well.
    Thus $\sigma(\mathcal S) = \mathcal M^B$ --- that is,
    if $B\in\sigma(\mathcal T)$ and $A\in\sigma(\mathcal S) = \mathcal M^B$,
    then $A\times B\in\sigma(\mathcal S\times\mathcal T)$,
    so $\sigma(\mathcal S\times\mathcal T)$ contains
    $\sigma(\mathcal S)\times\sigma(\mathcal T)$,
    hence also contains the generated $\sigma$-algebra
    $\sigma(\mathcal S)\otimes\sigma(\mathcal T)$.
    \\

    (b) Let $X$ be nonempty, $\mathcal S = 2^X$, and $\mathcal T = \{\varnothing\}$.
    Then $\sigma(\mathcal S) = 2^X$ and
    $\sigma(\mathcal T) = \{\varnothing, Y\}$, so
    \begin{equation*}
        \sigma(\mathcal S)\otimes\sigma(\mathcal T)
	= \{A\times Y : A\subseteq X\},
    \end{equation*}
    but $\mathcal S\times\mathcal T
    = \{A\times B : A\in\mathcal S, B\in\mathcal T\} = \{\varnothing\}$,
    so $\sigma(\mathcal S\times\mathcal T) = \{\varnothing\}$.
\end{soln}

\pagebreak

\begin{problem}
    Let $(X,\mathcal S)$ and $(Y,\mathcal T)$ be measurable spaces
    and let $f\colon X\to Y$ be a $(\mathcal T,\mathcal S)$-measurable function.
    \begin{enumerate}[(a)]
	\ii For $E\subseteq X\times Y$, let $E_f \coloneq \{x\in X: (x,f(x))\in E\}$.
	Show that $E_f\in\mathcal S$ for all $E\in\mathcal S\otimes\mathcal T$.
	\ii Let $A\subseteq X$. Show that $A\in\mathcal S$ if
	$\gr(f|_A) \coloneq \{(x,f(x)) : x\in A\}\in\mathcal S\otimes\mathcal T$.
	\ii Suppose $Y = \RR$ and $\mathcal T$ contains the Borel $\sigma$-algebra.
	Show that $\gr(f|_A)\in\mathcal S\otimes\mathcal T$ for all $A\in\mathcal S$.
    \end{enumerate}
\end{problem}

\begin{soln}
    (a) Let $\mathcal M$ be the collection of subsets $E\subseteq X\times Y$
    for which $E_f\in\mathcal S$; we want to show that
    $\mathcal S\otimes\mathcal T\subseteq\mathcal M$.
    If $A\in\mathcal S$ and $B\in\mathcal T$,
    then $(A\times B)_f = A\cap f\inv(B)$ because
    $(x,f(x))\in A\times B$ if and only if
    $x\in A$ and $f(x)\in B$, or equivalently $x\in f\inv(B)$.
    By assumption $f\inv(B)$ is in $\mathcal S$, so $(A\times B)_f\in\mathcal S$,
    hence $A\times B\in\mathcal M$.
    Furthermore, $\mathcal M$ is a $\sigma$-algebra because
    \begin{enumerate}[(1)]
	\ii $\varnothing\in\mathcal M$,
	since $\varnothing_f = \varnothing\in\mathcal S$;
	\ii $\mathcal M$ is closed under countable unions,
	since if $E_1,E_2,\ldots\in\mathcal M$, then
	\begin{equation*}
	    \left(\bigcup_{n=1}^\infty E_n\right)_f
	    = \{x\in X : (x,f(x))\in E_n\text{ for some }n\in\NN\}
	    = \bigcup_{n=1}^\infty (E_n)_f\in\mathcal S
	\end{equation*}
	so $\bigcup_{n=1}^\infty E_n\in\mathcal M$; and
	\ii $\mathcal M$ is closed under complements,
	since if $E\in\mathcal M$, then
	\begin{equation*}
	    (X\setminus E)_f
	    = \{x\in X : (x,f(x))\in X\setminus E\}
	    = X\setminus E_f\in\mathcal S
	\end{equation*}
	so $X\setminus E\in\mathcal M$.
    \end{enumerate}
    Since $\mathcal M$ is a $\sigma$-algebra containing
    all rectangles $A\times B$ where $A\in\mathcal S$ and $B\in\mathcal T$,
    $\mathcal M$ contains the $\sigma$-algebra generated
    by these rectangles, namely $\mathcal S\otimes\mathcal T$.
    \\

    (b) In the notation of part (a),
    \begin{equation*}
	(\gr(f|_A))_f = \{x\in X : (x,f(x))\in\gr(f|_A)\} = A,
    \end{equation*}
    so part (a) implies that if $\gr(f|_A)\in\mathcal S\otimes\mathcal T$
    then $A\in\mathcal S$.
    \\

    (c) Let $\QQ = \{q_1,q_2,\ldots\}$ be an enumeration of the rationals.
    For any $n\in\NN$, the set
    \begin{equation*}
	X_n = \bigcup_{k=1}^\infty f\inv((q_k-1/n, q_k+1/n))\times
	(q_k-1/n, q_k+1/n)
    \end{equation*}
    is in $\mathcal S\otimes\mathcal T$ because $f$ is $\mathcal S$-measurable.
    Also, $X_n$ contains the entire graph $\gr(f)$
    because the set of intervals $\{(q_k-1/n, q_k+1/n) : k,n\in\NN\}$
    covers $\RR$, as every real number is within $1/n$ of some rational number.

    I claim that $\bigcap_{n=1}^\infty X_n = \gr(f)$.
    To show this, let $(x,y)\in X\times\RR$ be a point outside $\gr(f)$,
    i.e. $y\ne f(x)$. Then, pick $n$ large enough that
    $2/n < \abs{f(x) - y}$. For any rational number $q$ such that
    $\abs{f(x) - q} < 1/n$, our choice of $n$ implies that
    $f(x)$ is in $(q-1/n,q+1/n)$ but $y$ isn't.
    Thus $(x,y)$ is not in $X_n$, meaning the intersection of the $X_n$
    does not contain any points outside $\gr(f)$.

    In summary, $\gr(f) = \bigcap_{n=1}^\infty X_n\in\mathcal S\otimes\mathcal T$,
    so $\gr(f|_A) = \gr(f)\cap A\times\RR$ is in $\mathcal S\otimes \mathcal T$
    for all $A\in\mathcal S$.
\end{soln}

\pagebreak 
\begin{problem}
    \begin{enumerate}[(a)]
	\ii Let $\mathcal M\subseteq 2^\RR$ denote
	the collection of Lebesgue measurable subsets.
	Give an example of a subset $E\subseteq\RR\times\RR$ such that
	all horizontal and vertical sections $E_x,E^y\subseteq\RR$
	are Lebesgue measurable, but $E\notin\mathcal M\otimes\mathcal M$.
	\ii Let $\lambda$ denote the Lebesgue measure on $\mathcal M$.
	Show that the measure $\lambda\times\lambda$
	on $\mathcal M\otimes\mathcal M$ is not complete.
    \end{enumerate}
\end{problem}

\begin{soln}
    (a) Let $V\in 2^\RR\setminus\mathcal M$ be
    a non-Lebesgue-measurable subset of $\RR$.
    Then $E = \{(x,x) : x\in V\}$ is not in $\mathcal M\otimes\mathcal M$
    by the contrapositive of 2(b), where $f(x) = x$ and $A = V$.
    On the other hand, given $x,y\in\RR$,
    \begin{equation*}
        E_x = \begin{cases}
	    \{x\} & \text{if } x\in E\\
	    \varnothing & \text{if }x\notin E
	\end{cases}\qquad\text{ and }\qquad
        E^y = \begin{cases}
	    \{y\} & \text{if } y\in E\\
	    \varnothing & \text{if }y\notin E
        \end{cases}
    \end{equation*}
    which are each Lebesgue measurable.
    \\

    (b) Let $A = \{(x,x) : x\in\RR\}\subseteq\RR^2$.
    This is in $\mathcal M\otimes\mathcal M$ by 2(c),
    since $\mathcal M$ contains the Borel $\sigma$-algebra,
    and since $A_x = \{x\}$ for all $x\in\RR$,
    \begin{equation*}
        \lambda\times\lambda(A) =
	\int_\RR\left(\int_\RR \ind_A(x,y)\dd\lambda(y)\right)\dd\lambda(x)
	= \int_\RR \lambda(A_x)\dd\lambda = 0.
    \end{equation*}
    If $E$ is as in part (a), then $E\subseteq A$
    is a subset of a measure zero set that is not in $\mathcal M\otimes\mathcal M$.
\end{soln}

\begin{problem}
    [Axler 5A \#10]
    Suppose $(X,\mathcal S,\mu)$ and $(Y,\mathcal T,\nu)$
    are $\sigma$-finite measure spaces. Prove that if
    $\omega$ is a measure on $\mathcal S\otimes\mathcal T$
    such that $\omega(A\times B) = \mu(A)\nu(B)$
    for all $A\in\mathcal S$ and $B\in\mathcal T$,
    then $\omega = \mu\times\nu$.
\end{problem}

\begin{soln}
    Let
    \begin{equation*}
	\mathcal M = \{ E\in\mathcal S\otimes\mathcal T :
	\omega(E) = \mu\times\nu(E)\}
    \end{equation*}
    and let $\mathcal A$ be the collection of
    unions of finitely many (pairwise disjoint) sets in $\mathcal S\times\mathcal T$;
    this is an algebra on $X\times Y$ by Axler's 5.13.
    We show that $\mathcal M$ is a monotone class containing
    $\mathcal A$, which implies the result
    because $\mathcal M$ then contains monotone class containing $\mathcal A$,
    which is equal to $\sigma(\mathcal A) = \mathcal S\otimes\mathcal T$
    by Axler's 5.17.

    First, we show that $\mathcal A\subseteq\mathcal M$.
    Let $E\in\mathcal A$ be expressed as the union of rectangles
    \begin{equation*}
	E = \bigcup_{k=1}^m (A_k\times B_k),\quad
	A_1,\dots,A_m\in\mathcal S,\; A_i\cap A_j=\varnothing\text{ if }i\ne j,\quad
	B_1,\dots,B_m\in\mathcal T,\; B_i\cap B_j=\varnothing\text{ if }i\ne j.
    \end{equation*}
    Then, $E\in\mathcal M$ because
    \begin{equation*}
	\omega(E) = \sum_{k=1}^m \omega(A_k\times B_k)
	= \sum_{k=1}^m \mu(A_k)\nu(B_k)
	= \mu\times\nu(E).
    \end{equation*}

    If $E_1\subseteq E_2\subseteq\cdots$ is
    an increasing sequence of sets in $\mathcal M$, then
    \begin{equation*}
	\omega\left(\bigcup_{n=1}^\infty E_n\right)
	= \lim_{n\to\infty} \omega(E_n)
	= \lim_{n\to\infty} \mu\times\nu(E_n)
	= \mu\times\nu\left(\bigcup_{n=1}^\infty E_n\right),
    \end{equation*}
    so $\bigcup_{n=1}^\infty E_n\in\mathcal M$ too,
    i.e. $\mathcal M$ is closed under increasing countable unions.

    Now we show closure under decreasing countable intersections.
    Since $X$ and $Y$ are $\sigma$-finite, we can write
    \begin{align*}
	X &= \bigcup_{k=1}^\infty X_k,\;
	X_k\in\mathcal S,\;\mu(X_k) < \infty,\; X_k\subseteq X_{k+1}
	\text{ for all }k\in\NN,\\
	Y &= \bigcup_{k=1}^\infty Y_k,\;
	Y_k\in\mathcal T,\;\nu(Y_k) < \infty,\; Y_k\subseteq Y_{k+1}
	\text{ for all } k\in\NN.
    \end{align*}
    Then, the sequence $X_k\times Y_k$ is an increasing sequence
    of sets of finite measure (with respect to both $\omega$ and $\mu\times\nu$)
    whose union is $X\times Y$. The measurable space
    $(X_k\times Y_k,(\mathcal S|_{X_k})\otimes(\mathcal T|_{Y_k}))$,
    can be endowed with the restricted finite measures
    $\omega = \omega|_{X_k\times Y_k}$
    and $\mu\times\nu = (\mu\times\nu)|_{X_k\times Y_k}$.
    Let $\mathcal M_k = \{E\in(\mathcal S|_{X_k})\otimes(\mathcal T|_{Y_k})
    : \omega(E) = \mu\times\nu(E)\}$. Then, $\mathcal M_k$
    has the properties of $\mathcal M$ established above:
    it contains the algebra generated by finite unions of rectangles
    $A\times B$ for $A\in\mathcal S|_{X_k}$ and $B\in\mathcal T|_{Y_k}$,
    and it is closed under increasing countable unions.
    Furthermore, for any sequence
    $E_1,E_2,\ldots\in\mathcal M_k$ such that $E_1\supseteq E_2\supseteq\cdots$,
    \begin{equation*}
	\mu\times\nu\left(\bigcap_{n=1}^\infty E_n\right)
	= \lim_{n\to\infty} \mu\times\nu(E_n)
	= \lim_{n\to\infty} \omega(E_n)
	= \omega\left(\bigcap_{n=1}^\infty E_n\right).
    \end{equation*}
    Thus $\mathcal M_k$ is a monotone class, meaning it contains
    (hence equals) $(\mathcal S\otimes\mathcal T)|_{X_k\times Y_k}$.

    Finally, given $E\in\mathcal M$, we have that
    $E\cap (X_k\times Y_k)$ is in $(\mathcal S\otimes\mathcal T)|_{X_k\times Y_k}$,
    which we mentioned in class is equal to
    $(\mathcal S\otimes\mathcal T)|_{X_k\times Y_k}$, so
    \begin{align*}
	\omega\left(\bigcap_{n=1}^\infty E_n\right)
	&= \omega\left(\bigcup_{k=1}^\infty\bigcap_{n=1}^\infty
	E_n\cap(X_k\times Y_k)\right)
	= \lim_{k\to\infty}\lim_{n\to\infty}
	\omega(E_n\cap(X_k\times Y_k))\\
	&= \lim_{k\to\infty}\lim_{n\to\infty}
	\mu\times\nu(E_n\cap(X_k\times Y_k))
	= \mu\times\nu\left(\bigcup_{k=1}^\infty\bigcap_{n=1}^\infty
	E_n\cap(X_k\times Y_k)\right)
	= \mu\times\nu\left(\bigcap_{n=1}^\infty E_n\right)
    \end{align*}
    so $\bigcap_{n=1}^\infty E_n\in\mathcal M$.
    This finishes the proof that $\mathcal M$ is a monotone class,
    so it contains (hence equals) $\mathcal S\otimes\mathcal T$.
\end{soln}

\begin{problem}
    Let $([0,1],\mathcal M,\lambda)$ be the usual Lebesgue measure space,
    $\mu$ be the counting measure on $2^{[0,1]}$, and
    \begin{equation*}
	E = \{(x,x) : x\in[0,1]\}\subseteq[0,1]^2.
    \end{equation*}
    Show that
    \begin{equation*}
        E\in\mathcal M\otimes 2^{[0,1]},\qquad
	\mu(E_x) = 1\text{ for all }x\in[0,1],\qquad
	\lambda(E^y) = 0\text{ for all }y\in[0,1].
    \end{equation*}
    Why doesn't this contradict Axler's equation 5.29?
\end{problem}

\begin{soln}
    $E\in\mathcal M\otimes 2^{[0,1]}$ by problem 2(b), with $f(x) = x$
    being the identity function and $A = [0,1]$.
    For any $x\in[0,1]$, $E_x = E^x = \{x\}$, so
    $\mu(E_x) = 1$ and $\lambda(E^y) = 0$.
    This does not contradict Axler's 5.29 because
    $\mu$ is not $\sigma$-finite: if $A_1,A_2,\ldots\subseteq[0,1]$
    are subsets with $\mu(A_k)$ finite for all $k\in\NN$, then
    each $A_k$ is a finite set, so $\bigcup_{n=1}^\infty A_n$
    is countable, hence not equal to the uncountable set $[0,1]$.
\end{soln}

\begin{problem}
    Let $\lambda$ be the Lebesgue measure on $[0,1]$.
    For each $n\in\NN$, let $f_n\colon[0,1]\to\RR$ be
    a continuous function such that
    \begin{equation*}
	f_n(x) = 0\text{ for all }x\notin[2^{-n},2^{-n+1}],\qquad
	\int_0^1f_n\dd x = 1.
    \end{equation*}
    \begin{enumerate}[(a)]
        \ii Show that the sum
	\begin{equation*}
	    f(x,y) = \sum_{n=1}^\infty (f_n(x) - f_{n+1}(x))f_n(y)
	\end{equation*}
	converges for all $(x,y)\in[0,1]^2$ and the function
	$f\colon [0,1]^2-\{(0,0)\}\to\RR$ is continuous.
	\ii Show that
	\begin{equation*}
	    \int_{[0,1]}\left(\int_{[0,1]}f(x,y)\dd\lambda(y)\right)\dd\lambda(x) = 1
	    \qquad\text{and}\qquad
	    \int_{[0,1]}\left(\int_{[0,1]}f(x,y)\dd\lambda(x)\right)\dd\lambda(y) = 0.
	\end{equation*}
	Why doesn't this contradict Tonelli's/Fubini's theorems?
    \end{enumerate}
\end{problem}

\begin{soln}
    (a) Let $g_n\colon[0,1]^2\to\RR$ be the function
    $g_n(x,y) = (f_n(x) - f_{n+1}(x))f_n(y)$, which is continuous on $[0,1]^2$
    and zero outside the rectangle
    $I_n = [2^{-n-1}, 2^{-n+1}]\times[2^{-n},2^{-n+1}]$.
    Any open set $U\subseteq[0,1]^2$
    containing $(0,0)$ contains cofinitely many of the $I_n$
    because all the endpoints of the intervals being multiplied to make $I_n$
    converge to zero as $n$ grows arbitrarily large.

    Given $(x_0,y_0)\ne(0,0)$ in $[0,1]$, let $U$ be an open set containing $(x_0,y_0)$
    whose complement contains a neighborhood of $(0,0)$;
    concretely, one can pick $U = [0,1]^2\setminus[0,1/m]^2$ for
    $m$ such that $1/m < \max\{x_0,y_0\}$.
    Then, $U$ intersects finitely many of the $I_n$,
    so the sum $f(x,y) = \sum_{n=1}^\infty g_n(x,y)$ is really a finite sum
    over $U$, as only finitely many of the addends are nonzero on $U$.
    This means the restriction of $f$ to $U$ is a continuous function,
    and since $U$ is an open subset of $[0,1]^2$,
    this implies continuity of $f$ at $(x_0,y_0)\in U$.
    \\

    (b) For a given $x_0\in(0,1]$, if $2^{-n+1} < x_0$ then
    $f_n(x_0) = f_{n+1}(x_0) = 0$, so the section $(g_n)_{x_0}$ is zero everywhere,
    meaning $f(x_0,y) = \sum_{n=1}^N g_n(x_0,y)$ where $N\in\NN$ is
    such that $2^{-N+1} < x_0$. Then,
    \begin{align*}
	\int_{[0,1]} f(x_0,y)\dd\lambda(y)
	&= \sum_{n=1}^N \left((f_n(x_0) - f_{n+1}(x_0))
	\int_{[0,1]} f_n(y)\dd\lambda(y)\right)\\
	&= \sum_{n=1}^N f_n(x_0) - f_{n+1}(x_0) = f_1(x_0) - f_{N+1}(x_0) = f_1(x_0).
    \end{align*}
    Finally, $(0,1]$ and $[0,1]$ differ by a set of measure zero, so
    \begin{equation*}
	\int_{[0,1]}\left(\int_{[0,1]}f(x,y)\dd\lambda(y)\right)\dd\lambda(x)
	= \int_{(0,1]} \left(\int_{[0,1]}f(x,y)\dd\lambda(y)\right)\dd\lambda(x)
	= \int_{(0,1]} f_1\dd\lambda = 1.
    \end{equation*}

    To compute the other iterated integral we use the same procedure:
    for $y_0\in(0,1]$, the section $(g_n)^{y_0}$ is zero for $x\in[0,1]$
    if $2^{-n+1} < y_0$, so picking $N$ such that $2^{-N+1} < y_0$,
    \begin{equation*}
        \int_{[0,1]} f(x,y_0)\dd\lambda(x)
	= \sum_{n=1}^N \left(f_n(y_0)\cdot
	\int_{[0,1]} f_n(x) - f_{n+1}(x)\dd\lambda(x)\right)
	= \sum_{n=1}^n f_n(y_0)(1-1) = 0,
    \end{equation*}
    hence
    \begin{equation*}
	\int_{[0,1]}\left(\int_{[0,1]}f(x,y)\dd\lambda(y)\right)\dd\lambda(x)
	= \int_{(0,1]}\left(\int_{[0,1]}f(x,y)\dd\lambda(y)\right)\dd\lambda(x)
	= \int_{(0,1]} 0\dd\lambda = 0.
    \end{equation*}

    Now we check that we have not encountered a contradiction
    to Tonelli's/Fubini's theorems.
    Note that the functions $a_n,b_n\colon[0,1]^2\to\RR$ given by
    $a_n(x,y) = f_n(x)f_n(y)$ and $b_n(x,y) = f_{n+1}(x)f_n(y)$
    are zero outside the rectangles $I_n = (2^{-n},2^{-n+1}]\times(2^{-n},2^{-n+1}]$
    and $J_n = (2^{-n-1},2^{-n}]\times(2^{-n},2^{-n+1}]$ respectively
    (where the intervals are half-open because continuity guarantees
    that $f_n$ is zero on the closure of $[0,1]\setminus[2^{-n},2^{-n+1}]$,
    hence zero on $[0,2^{-n}]$).
    These rectangles are pairwise disjoint, so $f|_{I_n} = a_n$ and $f|_{J_n} = -b_n$.
    By the integral condition on $f_n$ there exists a set of positive measure $A_n$
    on which $f_n$ is positive for each $n$; therefore
    $-b_n(x,y) = -f_n(x)f_{n+1}(y)$ is negative on $A_n\times A_{n+1}$,
    so $f$ is negative on $A_n\times A_{n+1}$,
    meaning Tonelli's theorem doesn't apply here
    (even if we only require almost-everywhere nonnegativity).
    Furthermore,
    \begin{align*}
        \int_{[0,1]^2}\abs f
	\dd(\lambda\times\lambda)
	&= \int_{[0,1]^2}\left(
	\sum_{n=1}^\infty \abs{a_n} + \abs{b_n}\right)
	\dd(\lambda\times\lambda)
	= \sum_{n=1}^\infty\left(\int_{[0,1]^2}
	\abs{a_n} + \abs{b_n}
	\dd(\lambda\times\lambda)\right)\\
	&= \sum_{n=1}^\infty\int_{[0,1]}\int_{[0,1]}
	\abs{f_n(y)}(\abs{f_n(x)} + \abs{f_{n+1}(x)})
	\dd\lambda(x)\dd\lambda(y)\\
	&\ge \sum_{n=1}^\infty\int_{[0,1]}2\abs{f_n(y)}
	\dd\lambda(x)\dd\lambda(y) \ge \sum_{n=1}^\infty 2 = \infty
    \end{align*}
    where the first equality is by disjointness of the $I_n$ and $J_n$;
    the second equality is by monotone convergence;
    and the third is by Tonelli's theorem.
    Therefore, the hypothesis of Fubini's theorem does not hold here.
\end{soln}










\end{document}
